% !TeX root = ../algebra-1.tex
\chapter{Linear maps and coordinates}\label{chap:algebra-linear-maps}

\section{Matrices as linear maps}

An $m\times n$ matrix defines a function
\[
 T_A:F^n\longrightarrow F^m,\qquad T_A(v)=Av.
\]
Its input has $n$ coordinates and its output has $m$ coordinates.
For example, $\begin{pmatrix}-1&0\\0&1\end{pmatrix}$ reflects
$\mathbb R^2$ in the $y$-axis.

\begin{definition}
A map $T:F^n\to F^m$ is \emph{linear} if
\[
 T(u+v)=T(u)+T(v),\qquad T(av)=aT(v)
\]
for all $u,v\in F^n$ and $a\in F$.
Equivalently, $T(au+bv)=aT(u)+bT(v)$ for all $a,b,u,v$.
\end{definition}

Matrix multiplication shows that $T_A$ is linear. Repeated application of
linearity gives
$T(\sum_i a_iv_i)=\sum_i a_iT(v_i)$.
Every linear map sends zero to zero, since $T(0v)=0T(v)=0_m$.
A nonzero translation $v\mapsto v+c$ is therefore not linear.

\begin{proposition}[Composition]
If $B$ is $n\times p$ and $A$ is $m\times n$, then
\[
 T_A\circ T_B=T_{AB}.
\]
\end{proposition}
\begin{proof}
For $v\in F^p$, $T_A(T_B(v))=A(Bv)=(AB)v$.
\end{proof}

The rightmost map acts first. An invertible matrix represents a linear
bijection, and its inverse map is represented by the inverse matrix.

\begin{theorem}[Every linear map has a matrix; source Theorem 3.4]
For each linear $T:F^n\to F^m$, there is a unique $m\times n$ matrix $A$
such that $T(v)=Av$ for every $v$. Its columns are
\[
 \boxed{A=[\,T(e_1)\ \cdots\ T(e_n)\,].}
\]
\end{theorem}
\begin{proof}
Write $v=\sum_j v_je_j$. Linearity gives
$T(v)=\sum_j v_jT(e_j)$, exactly the column formula for $Av$.
Conversely, a representing matrix must have column $j$ equal to
$Ae_j=T(e_j)$, so it is unique.
\end{proof}

\begin{example}
The map $T(a,b,c)=a+b+c$ has matrix $(1\ 1\ 1)$.
The map $S(a,b)=(a,b,a,b)^{\mathsf T}$ has matrix
\[
 \begin{pmatrix}1&0\\0&1\\1&0\\0&1\end{pmatrix}.
\]
Finding images of the standard basis is often quicker than guessing
matrix entries.
\end{example}

\section{Kernel, image and invertibility of a map}

\begin{definition}
The \emph{kernel} and \emph{image} of a linear map are
\[
 \ker T=\{v:T(v)=0_m\},\qquad
 \im T=\{T(v):v\in F^n\}.
\]
If $A$ represents $T$, these are $N(A)$ and $C(A)$.
Its rank is $\dim(\im T)$.
\end{definition}

They are subspaces by the same closure proofs as for matrices.
Rank--nullity becomes
\[
 \dim(\ker T)+\dim(\im T)=n.
\]

\begin{example}
For $T(a,b,c)=a+b+c$, the kernel is the plane $a+b+c=0$,
with basis $(1,0,-1)^{\mathsf T},(0,1,-1)^{\mathsf T}$.
The image is all of $F$, since $T(t,0,0)=t$.
Thus the kernel dimension is two and the rank is one.

For $S(a,b)=(a,b,a,b)^{\mathsf T}$, the kernel is $\{0_2\}$ and
\[
 \im S=\spanop\left(
 (1,0,1,0)^{\mathsf T},(0,1,0,1)^{\mathsf T}\right).
\]
This map is injective, but is not surjective onto $F^4$.
\end{example}

\begin{proposition}
A linear map is injective if and only if its kernel is $\{0_n\}$.
It is surjective if and only if its image is the whole codomain.
\end{proposition}
\begin{proof}
If it is injective, $T(v)=0_m=T(0_n)$ implies $v=0_n$.
Conversely, if its kernel is trivial and $T(u)=T(v)$, then
$T(u-v)=0_m$ gives $u=v$.
The second assertion is the definition of surjectivity.
\end{proof}

Thus for $T:F^n\to F^m$, injectivity means rank $n$ and surjectivity
means rank $m$. Rank--nullity implies that an injective map requires
$n\leq m$ and a surjective map requires $n\geq m$.
When $m=n$, injectivity and surjectivity are equivalent, and both mean
that the representing matrix is invertible.

\section{Coordinates in an ordered basis}

\begin{definition}
For an ordered basis $\mathcal B=(b_1,\ldots,b_n)$ of $F^n$, the
\emph{coordinate vector} of $v$ is
\[
 [v]_{\mathcal B}=\begin{pmatrix}c_1\\\vdots\\c_n\end{pmatrix}
 \quad\text{where}\quad v=c_1b_1+\cdots+c_nb_n.
\]
\end{definition}

Let $P=[\,b_1\ \cdots\ b_n\,]$. Its columns form a basis, so $P$ is
invertible and
\[
 \boxed{v=P[v]_{\mathcal B},\qquad [v]_{\mathcal B}=P^{-1}v.}
\]
The vector $v$ is unchanged; these equations describe the same vector
using different coordinate systems.
The direction is worth remembering:
\[
 \mathcal B\text{-coordinates}
 \ \xrightarrow{\ P\ }\ \text{standard coordinates},
 \qquad
 \text{standard coordinates}
 \ \xrightarrow{\ P^{-1}\ }\ \mathcal B\text{-coordinates}.
\]

\begin{example}
Take $b_1=(1,1)^{\mathsf T}$ and $b_2=(1,-1)^{\mathsf T}$.
Then
\[
 P=\begin{pmatrix}1&1\\1&-1\end{pmatrix},\qquad
 P^{-1}=\frac12\begin{pmatrix}1&1\\1&-1\end{pmatrix}.
\]
For $v=(0,1)^{\mathsf T}$, $[v]_{\mathcal B}=(1/2,-1/2)^{\mathsf T}$:
\[
 \begin{pmatrix}0\\1\end{pmatrix}
 =\frac12\begin{pmatrix}1\\1\end{pmatrix}
  -\frac12\begin{pmatrix}1\\-1\end{pmatrix}.
\]
\end{example}

\section{Changing the matrix of a map}

\begin{theorem}[Change of basis; source Theorem 3.14]
Suppose $T:F^n\to F^m$ has standard matrix $A$.
Let $\mathcal B$ be an input basis with column matrix $P$, and
$\mathcal C$ an output basis with column matrix $Q$.
The matrix taking $\mathcal B$-coordinates of the input to
$\mathcal C$-coordinates of the output is
\[
 \boxed{[T]_{\mathcal C\leftarrow\mathcal B}=Q^{-1}AP.}
\]
\end{theorem}
\begin{proof}
If $x=[v]_{\mathcal B}$, then $v=Px$.
The standard output is $APx$, and its $\mathcal C$-coordinates are
$Q^{-1}APx$. Equivalently, column $j$ of the new matrix is
$[T(b_j)]_{\mathcal C}$.
\end{proof}

When $T:F^n\to F^n$ uses the same basis for its input and output, this
becomes $P^{-1}AP$.
Matrices related by $B=P^{-1}AP$ are called \emph{similar}.
Similarity describes one operator in different bases.
Row reduction instead multiplies on the left by elementary matrices;
it serves a different purpose and is not generally similarity.

\begin{example}[A basis adapted to the map]
Let $A=\begin{pmatrix}2&1\\1&2\end{pmatrix}$ and use the preceding basis.
Since
\[
 A\begin{pmatrix}1\\1\end{pmatrix}
 =3\begin{pmatrix}1\\1\end{pmatrix},\qquad
 A\begin{pmatrix}1\\-1\end{pmatrix}
 =\begin{pmatrix}1\\-1\end{pmatrix},
\]
we obtain
\[
 P^{-1}AP=\begin{pmatrix}3&0\\0&1\end{pmatrix}.
\]
In these coordinates the map stretches one coordinate by $3$ and leaves
the other unchanged.
\end{example}

\section{Diagonalisation and matrix powers}

If $D=\diag(d_1,\ldots,d_n)$, then
$D^k=\diag(d_1^k,\ldots,d_n^k)$ for positive integers $k$.
When $P^{-1}AP=D$, equivalently $A=PDP^{-1}$, cancellation of the middle
factors gives
\[
 \boxed{A^k=PD^kP^{-1}\quad(k\geq0).}
\]
Here $D^0=I_n$. Negative powers are also valid when all $d_i$ are
nonzero, so that $A$ is invertible.

\begin{example}
For $A=\begin{pmatrix}3&-1\\-1&3\end{pmatrix}$ and the same $P$,
\[
 P^{-1}AP=\diag(2,4).
\]
Therefore, for $k\geq0$,
\[
 A^k=\frac12
 \begin{pmatrix}
 2^k+4^k&2^k-4^k\\
 2^k-4^k&2^k+4^k
 \end{pmatrix}.
\]
The directions $(1,1)^{\mathsf T}$ and $(1,-1)^{\mathsf T}$ are
scaled by $2^k$ and $4^k$ respectively.
\end{example}

\subsection{Application: the Fibonacci sequence}

Define $f_0=0$, $f_1=1$ and $f_{k+2}=f_{k+1}+f_k$.
With
\[
 M=\begin{pmatrix}1&1\\1&0\end{pmatrix},
 \qquad
 \begin{pmatrix}f_{k+1}\\f_k\end{pmatrix}
 =M^k\begin{pmatrix}1\\0\end{pmatrix},
\]
the recurrence becomes a matrix power.
To find invariant directions, calculate
$M(\lambda,1)^{\mathsf T}=(\lambda+1,\lambda)^{\mathsf T}$.
This equals $\lambda(\lambda,1)^{\mathsf T}$ when
$\lambda^2=\lambda+1$.
The two roots are
\[
 \phi_+=\frac{1+\sqrt5}{2},\qquad
 \phi_-=\frac{1-\sqrt5}{2}.
\]
Set
\[
 P=\begin{pmatrix}\phi_+&\phi_-\\1&1\end{pmatrix},\qquad
 P^{-1}=\frac1{\sqrt5}
 \begin{pmatrix}1&-\phi_-\\-1&\phi_+\end{pmatrix}.
\]
The invariant-direction calculations give
$MP=P\diag(\phi_+,\phi_-)$.
Hence
\begin{align*}
 \begin{pmatrix}f_{k+1}\\f_k\end{pmatrix}
 &=P\begin{pmatrix}\phi_+^k&0\\0&\phi_-^k\end{pmatrix}
       \frac1{\sqrt5}\begin{pmatrix}1\\-1\end{pmatrix}\\
 &=\frac1{\sqrt5}
   \begin{pmatrix}\phi_+^{k+1}-\phi_-^{k+1}\\
                   \phi_+^k-\phi_-^k\end{pmatrix}.
\end{align*}
The second coordinate is \emph{Binet's formula}:
\[
 \boxed{f_k=\frac{\phi_+^k-\phi_-^k}{\sqrt5}.}
\]
The invariant directions used here are eigenvectors, the main topic of
the next chapter.
